\section{Verificação de rastreabilidade e coerência}
\begin{table}[!ht]\centering\small\renewcommand{\arraystretch}{1.5}
\begin{tabularx}{\linewidth}{|>{\raggedright\arraybackslash}p{4cm}|>{\raggedright\arraybackslash}X|}
\hline
\textbf{Ponto de verificação} & \textbf{Resposta do grupo} \\
\hline
Os requisitos foram considerados? & Todos os RF01--RF19 possuem dados ou relacionamentos associados no MER, no DER ou nas entidades associativas. \\
\hline
Os identificadores estão definidos? & Cada entidade possui chave própria; os relacionamentos muitos para muitos são representados por chaves compostas. \\
\hline
As regras de negócio estão registradas? & Moderação, segmentação, recebimento, parentesco, matrícula, grade acadêmica e mensagens possuem regras e dúvidas de validação. \\
\hline
Existem entidades sem uso rastreável? & REACAO\_AVISO e COMENTARIO estão ligados ao RF19; NOTA e FREQUENCIA ampliam o acompanhamento acadêmico. \\
\hline
O modelo mantém o problema de origem? & O MER concentra comunicação centralizada, moderação, segmentação, entrega, recebimento, atendimento e acompanhamento escolar. \\
\hline
Quais pontos precisam ser validados? & Categorias sujeitas à moderação; vínculo entre responsáveis e alunos; canais externos; confirmação de leitura; comentários; horários e avaliações acadêmicas. \\
\hline
\end{tabularx}
\end{table}

\begin{objectives}
\textbf{Conclusão.} O modelo conceitual derivado da Entrega 2 e das fichas de MER/DER representa
as entidades necessárias ao funcionamento do sistema. Os diagramas foram preservados e as regras
de negócio foram organizadas como respostas às atividades solicitadas pela professora.
Pontos dependentes da validação institucional permanecem explicitamente registrados.
\end{objectives}
